\section{BIJLAGE · Spiekbrief en resultaten}

\begin{missie}
Missie: alles op één blad. Gebruik het tijdens het oefenen, leer het voor de toets.
\end{missie}

\subsection*{Centrummaten}

\breedtes{22/50/28}
\begin{tabular}{lll}
\hline
Maat & Hoe? & Bij welke data? \\
\hline
gemiddelde & $\overline{x}=\dfrac{\text{som}}{n}$ (steekproef), $\mu=\dfrac{\text{som}}{N}$ (populatie) & numeriek \\
mediaan & middelste waarde van de geordende data (even aantal: gemiddelde van de twee middelste) & numeriek, categorisch geordend \\
modus & waarde die het vaakst voorkomt & alle soorten data \\
\hline
\end{tabular}

\subsection*{Spreidingsmaten}

\breedtes{25/45/30}
\begin{tabular}{lll}
\hline
Maat & Formule & Opmerking \\
\hline
variatiebreedte & $R=\text{max}-\text{min}$ & gevoelig voor uitschieters \\
interkwartielafstand & $\Delta Q=Q_3-Q_1$ & middelste 50 \%; weinig gevoelig voor uitschieters \\
variantie & $\sigma^2=\dfrac{\sum (x-\mu)^2}{N}$, $\;s^2=\dfrac{\sum (x-\overline{x})^2}{n-1}$ & eenheid in het kwadraat \\
standaardafwijking & $\sigma=\sqrt{\sigma^2}$, $\;s=\sqrt{s^2}$ & dezelfde eenheid als de data \\
\hline
\end{tabular}

\subsection*{Steekproef en populatie}

\breedtes{34/33/33}
\begin{tabular}{lcc}
\hline
 & steekproef & populatie \\
\hline
omvang & $n$ & $N$ \\
gemiddelde & $\overline{x}$ & $\mu$ \\
variantie & $s^2$ & $\sigma^2$ \\
standaardafwijking & $s$ & $\sigma$ \\
\hline
\end{tabular}

Representatief = getrouwe afspiegeling van de populatie (aselecte trekking). Selectiebias = vertekening door zelfselectie, gemakssteekproef, onderdekking, overlevers of niet-respons.

\subsection*{Samenhang en causaliteit}

Samenhang $\neq$ causaliteit. Verklaringen: $A\Rightarrow B$, $B\Rightarrow A$, verborgen variabele $C$, toeval. Experiment: testgroep, controlegroep, toevallige indeling, dezelfde omstandigheden.

\subsection*{De normale verdeling}

\begin{theorie}
$X\sim N(\mu,\sigma)$: klokvormig, symmetrisch rond $\mu$, totale oppervlakte 1, buigpunten op $\mu\pm\sigma$. Model bij data: $\mu\approx\overline{x}$ en $\sigma\approx s$.

68 – 95 – 99,7-regel: $\mu\pm\sigma$: 68 \% \quad $\mu\pm 2\sigma$: 95 \% \quad $\mu\pm 3\sigma$: 99,7 \%.

Kansen: $P(X<a)$ = oppervlakte links, $P(X>a)=1-P(X<a)$, $P(a<X<b)=P(X<b)-P(X<a)$. TI: normalcdf(onder, boven, $\mu$, $\sigma$). GeoGebra: Normaal($\mu$, $\sigma$, $a$, Waar). Excel: =NORM.VERD($a$; $\mu$; $\sigma$; WAAR).
\end{theorie}

\subsection*{Mijn resultaten}

\breedtes{14/46/20/20}
\begin{tabular}{llcc}
\hline
Level & Onderwerp & Score CHECK & Behaald? \\
\hline
0 & Startklaar & \_\_ / 8 & ☐ \\
1 & Centrummaten & \_\_ / 7 & ☐ \\
2 & Spreidingsmaten & \_\_ / 10 & ☐ \\
3 & Steekproef en representativiteit (LPD 11) & \_\_ / 8 & ☐ \\
4 & Samenhang en causaliteit (LPD 12) & \_\_ / 9 & ☐ \\
5 & Normale verdeling als model (LPD 13) & \_\_ / 8 & ☐ \\
6 & Kansen berekenen (LPD 14) & \_\_ / 6 & ☐ \\
7 & Eindmissie (LPD 1 – 4) & & ☐ \\
Eind & Eindtoets & \_\_ / 60 & ☐ \\
\hline
\end{tabular}
